% eksperymentalne tłumaczenie na włoski

%

\subsection{Okrąg wpisany i trzy dwusieczne} % lang-pl
\subsection{Cerchio inscritto e tre bisettrici} % lang-it

\begin{proposition}
    Trzy dwusieczne kątów trójkąta przecinają się w jednym punkcie, środku okręgu wpisanego w ten trójkąt. % lang-pl
    Le tre bisettrici degli angoli di un triangolo si incontrano in un unico punto, il centro del cerchio inscritto nel triangolo. % lang-it
\end{proposition}

Uzasadnienie jest kompletnie analogiczne do istnienia środku okręgu opisanego (tak zrobi Neugebauer \cite[s. 32]{neugebauer_2018}), ale można też podać fikuśny dowód dookoła: z twierdzenia o dwusiecznej i Cevy, jak Zetel \cite[s. 14]{zetel_2020}. % lang-pl
Hartshorne \cite[s. 16]{hartshorne_2000} podaje to w formie ćwiczenia ze wskazówką, by spojrzeć na (IV.4), z czym totalnie się zgadzamy. % lang-pl
La dimostrazione è del tutto analoga a quella dell'esistenza del centro del cerchio circoscritto (come fa Neugebauer \cite[p. 32]{neugebauer_2018}), ma si può anche fornire una dimostrazione più elegante: mediante il teorema della bisettrice e il teorema di Ceva, come in Zetel \cite[p. 14]{zetel_2020}. % lang-it
Hartshorne \cite[p. 16]{hartshorne_2000} lo presenta come esercizio con il suggerimento di guardare (IV.4), e siamo pienamente d'accordo. % lang-it

Po angielsku mamy krótkie określenia \emph{excenter, excircle, incenter, incircle}. % lang-pl
In inglese si usano i termini brevi \emph{excenter, excircle, incenter, incircle}. % lang-it

% TODO: długość dwusiecznej, Zetel s. 32

\begin{proposition}
    Niech $r$ i $R$ będą promieniami okręgów wpisanego i opisanego na trójkącie o kątach $\alpha, \beta, \gamma$. % lang-pl
    Wtedy % lang-pl
    Siano $r$ e $R$ i raggi del cerchio inscritto e circoscritto di un triangolo con angoli $\alpha, \beta, \gamma$. % lang-it
    Allora % lang-it
    \begin{equation}
        \frac r R = 4 \sin \frac \alpha 2 \sin \frac \beta 2 \sin \frac \gamma 2.
    \end{equation}
\end{proposition}

Trygonometryczne rozwiązanie opisze Guzicki \cite[s. 351]{guzicki_2021}. % lang-pl
Una dimostrazione trigonometrica è presentata da Guzicki \cite[p. 351]{guzicki_2021}. % lang-it

%